\clearpage
\section[Parte N – Rev. 1.15 – Quarzo HSE Y1 (S-04)]{Parte N – Rev. 1.15 – Quarzo HSE Y1 (S-04)}

Punto S-04 della verifica \code{sch\_\allowbreak{}check\_\allowbreak{}2010.\allowbreak{}pdf}: il valore di Y1 non indicava frequenza e capacità di carico. Il circuito era già corretto: è stata corretta solo la documentazione sullo schema.

\subsection[27. S-04 – Quarzo Y1 dell'STM32L562]{27. S-04 – Quarzo Y1 dell'STM32L562 \tg{tagdone}{chiuso v1.15}}

{\small\arrayrulecolor{rulec}\setlength{\tabcolsep}{4pt}\renewcommand{\arraystretch}{1.1}
\begin{xltabular}{\linewidth}{>{\hsize=0.4977\hsize\raggedright\arraybackslash}X>{\hsize=2.2458\hsize\raggedright\arraybackslash}X>{\hsize=0.2565\hsize\raggedright\arraybackslash}X}
\toprule
\textbf{Controllo} & \textbf{Calcolo /\allowbreak{} dato} & \textbf{Esito} \\
\midrule
\endhead
\midrule[0.15pt]\multicolumn{3}{r}{\footnotesize\itshape continua sulla pagina seguente\ldots}\\\bottomrule\endfoot
\bottomrule\endlastfoot
Frequenza & 24 MHz su PH0/\allowbreak{}PH1 (HSE). L'HSE dell'STM32L562 accetta 4–48 MHz & OK \\\midrule[0.15pt]
C1, C2 & C = 2·(C\textsubscript{L} − C\textsubscript{s}) = 2·(10 − 5) = 10 pF, con C\textsubscript{s} ≈ 5 pF stimata. Montati 10 pF C0G 0402 (C76946). Nota col calcolo già presente sul foglio & OK \\\midrule[0.15pt]
R1 & 0 Ω, resistenza serie per limitare il drive level, aumentabile sul prototipo & OK \\\midrule[0.15pt]
Margine d'innesco & g\textsubscript{m,crit} = 4·ESR·(2πf)²·(C0 + C\textsubscript{L})² = 4·50·(2π·24 MHz)²·(3 + 10 pF)² ≈ 0.77 mA/\allowbreak{}V, con C0 ≈ 3 pF stimata (da confermare sul datasheet YXC). Limite STM32L5: 1.5 mA/\allowbreak{}V. Margine circa 2 & OK \\
\end{xltabular}}

\textbf{Per il firmware:} nel progetto STM32 impostare \code{HSE\_\allowbreak{}VALUE = 24000000}. Con un valore diverso UART, USB e timer partono alla frequenza sbagliata.

\textbf{Verifica.} Netlist identica alla rev. 1.14. ERC identico (138, nessuna nuova).
